\documentclass{IEEEtran}
\usepackage{cite}
\usepackage{amsmath,amssymb,amsfonts}
\usepackage{graphicx}
\usepackage[export]{adjustbox}
\usepackage{subcaption}
\captionsetup{font=small}
\usepackage{footmisc}
\usepackage{comment}
\usepackage{textcomp,nicefrac}
\def\BibTeX{{\rm B\kern-.05em{\sc i\kern-.025em b}\kern-.08em
T\kern-.1667em\lower.7ex\hbox{E}\kern-.125emX}}
\usepackage{dblfloatfix}

%\markboth{IEEE TRANSACTIONS ON NUCLEAR SCIENCE, VOL. XX, NO. XX, XXXX
%2025}
%{Author \MakeLowercase{\textit{et al.}}: Preparation of Papers for
%Review by the \textsc{IEEE Transactions on Nuclear 
%Science} \newline (Jan 2025)}

\input jidef
\input beginend
\input alphabet
\def\sym{^\text{sym}}
\def\ImDom{_\text{I}}
\def\CountDom{_\text{C}}
\def\CI{_\text{CI}}
\def\CountDomMC{_{\text{C-MC}}}
\def\ImDomMC{_{\text{I-MC}}}

\usepackage{color}
\usepackage{soul}
\usepackage{xspace}
\definecolor{violet}{rgb}{.56,0,1}
\definecolor{darkgreen}{rgb}{0,.4,0}
\definecolor{darkred}{rgb}{.6,0,0}
\definecolor{orange}{rgb}{1.,0.5,0}
%% MACROS pour identifier les auteurs 
\newcommand{\JI}[1]{{\color{darkgreen}{#1}}\xspace}
\newcommand{\SB}[1]{{\color{violet}{#1}}\xspace}
\newcommand{\SST}[1]{{\color{orange}{#1}}\xspace}

\newlength{\imagewidth}
\setlength{\imagewidth}{0.145\textwidth}
\usepackage{array}

\begin{document}
%\title{Learning Image Denoising from Single Scans for PET reconstruction}
%\title{Self-supervised PET Reconstruction Based on the Noise2Noise Principle}
\title{Self-supervised PET Denoising Based on the Noise2Noise Principle}
\author{Bouchez-Delotte Sacha$^{1}$, Jérôme Idier$^{1}$, Simon Stute$^{2}$, Clément Huneau$^{1}$

\thanks{$^{1}$Nantes Univ., École Centrale Nantes, CNRS, LS2N, %F-44300
Nantes, France
      }%
\thanks{$^{2}$CHU Nantes, Service de Médecine Nucléaire, CRCI$^2$NA, INSERM UMR 1307, Nantes University,%F-44000
Nantes, France
      }%
}

\maketitle
\begin{abstract}
PET reconstruction aims to transform raw sinograms data into images using analytical, iterative or deep learning approaches. In addition to the ill-posedness of the problem, inherent Poisson noise in sinograms makes the problem more difficult, especially in low-dose contexts.

Using deep learning for image denoising and reconstruction has shown promising performances. However, supervised training requires large sets of paired data with corrupted images and their associated ground truth, which is hard or even impossible to acquire in the medical field, especially in emission tomography.

In this work, we adapted Noise2Noise, a self-supervised framework, to PET imaging in order to learn how to denoise images from noisy measurements only.
%This method is shown to be equivalent to its supervised counterpart under some statistical assumptions that can be met with Poisson statistics. Indeed, sinograms can be split into two (or more) independent subsets of the original sinograms with higher signal-to-noise ratio using binomial sampling strategies.
%Just as any denoising estimator, the latter may be biased toward smoothness, which is obvious with PET imaging due to its low-resolution nature. In order to alleviate this issue and stabilize training, we also introduce measurement consistency terms in learning objectives.

%Another interest of our method is that it can either be implemented in the count domain using sinograms or in the image domain using reconstructed images. The MSE loss can be used in both applications, but count-domain training enables us to use more appropriate loss functions such as Poisson likelihood or Kullback-Leibler divergence.
\end{abstract}

\begin{IEEEkeywords}
Deep Learning, Denoising, PET Reconstruction, Noise2Noise, Self-supervised
\end{IEEEkeywords}

\section{Introduction}
\label{sec:introduction}
\IEEEPARstart{P}{ositron} emission tomography (PET) image reconstruction aims at retrieving an activity distribution $x$ from noisy measurements $y$ linked by the following forward model:
%\IEEEPARstart{P}{ositron} emission tomography (PET) is a medical imaging technique that requires the injection of a radiotracer into the patient. The aim of the reconstruction problem is to retrieve the radiotracer distribution within the patient tissues and organs from noisy sinogram data, i.e.\ retrieving $x$ from $y$ in the following foward model: % SIMON: on peut abbréger pour aller plus vite ici
%
%\begin{equation}
%    y = \Pc(Ax + b) + \epsilon % SIMON: there is no epsilon here
%\end{equation}
\begin{equation}
    y = \Pc(Ax + b)
\end{equation}
%where $A$ stands for a projection operator, $b$ contains systematic scatter and random contributions, and $\Pc(\lambda)$ stands for componentwise Poisson noise.
where $A$ is the projection operator modelling the PET acquisition system, $b$ contains background scatter and random contributions, and $\Pc(\lambda)$ stands for componentwise Poisson noise.
The filtered back projection (FBP) and the Maximum Likelihood Expectation Maximization (MLEM) algorithms respectively correspond to the most common analytical and algebraic non-regularized reconstruction techniques in PET.
%Most common non-regularized approaches include the filtered back projection (FBP) and the converged version of the  Maximum Likelihood Expectation Maximization (MLEM) algorithm.
Both are known to give noisy
results since the problem is ill-posed and the data is noisy. 

%Traditional strategies to alleviate this issue consists in introducing 
Regularization is necessary to tackle the ill-posed nature of the problem, stabilize the reconstruction process and reduce the noise level. For example, a classical solution is to maximize a penalized log-likelihood criterion involving a suited prior for PET such as the Relative Difference Prior (RDP)~\cite{Nuyts02}.

Deep learning-based methods constitute another possibility to compensate for the information deficiency of the data.
In the case of PET,
supervised deep learning-based denoising %for PET images is also an active research field. It
can be either performed in the count domain, where we learn to denoise sinograms and reconstruct the image from denoised sinogram at inference, or in the image domain, where we learn to denoise under-regularized reconstructed images as a post-processing step~\cite{Hashimoto2024}.

Supervised methods require large sets of paired data with corrupted images and their associated ground truth, which is hard or even impossible to acquire in the medical field, especially in emission tomography. This justifies the growing interest received by non-supervised methods.
For example, the unsupervised Deep Image Prior (DIP) framework -- where the only prior used for the denoising task is the image itself -- has been considered in PET imaging~\cite{Gong2019}.
Here, we rather focus on Noise2Noise (N2N), a self-supervised scheme that
%The self-supervised Noise2Noise training scheme
alleviates the lack of ground truth by relying on pairs of noisy realizations of objects~\cite{Lehtinen18}.
Given a set of such pairs $(y^1,y^2)=(y_i^1, y_i^2)_{n=1\ldotsv N}$, we train a neural network to map each noisy observation to its counterpart, based on a symmetrized loss function $\Lc\sym(y^1, y^2, \theta)=\Lc(y^1, y^2, \theta)+\Lc(y^2, y^1, \theta)$ with
$
\Lc(y^1, y^2, \theta) = \frac{1}{2} \sum_i%{n=1}^N
d(f_\theta(y_i^1), y_i^2),
$
$\theta$ being the parameters of a UNet model $f_\theta$ and $d$ a distance measure such as the Mean Squared Error (MSE).
Under the assumption of statistical independence between $y_i^1$ and $y_i^2$ and of centered noise, N2N becomes equivalent to its supervised counterpart when $N \to \infty$~\cite{Lehtinen18}.
To our best knowledge, the N2N principle has been only applied twice in the context of PET imaging. In~\cite{Kang2021}, independent Poisson noise realizations are applied to projected synthetic phantoms to get pairs, which is not feasible in a real clinical context since noise-free sinograms are not available in this case.
On the other hand,~\cite{Yie2020} learns to denoise 3D local patches from splitted list-mode real data, without assuming that the method may require more global context to learn the statistics.
%\SST{J'ai lu la reference Yie2020, ils ne proposent pas d'acquérir deux fois, ils disent que l'hypothèse derrière N2N revient à acquérir deux fois la même situation. Ils n'utilisent que des données réelles. A partir d'un listmode de 300s, ils fabriquent 100 list-mode de 10s avec un overlap de 7s entre chaque. Ensuite ils fabriquent des jeux de 20s en sommant 4 listmode de 10s sans overlap. Puis ils utilisent ces jeux pour l'entraitement sans dire comment ils appairent. La limitation dans cette méthode c'est que la distribution d'activité change entre deux splits et il y a l'effet de la décroissance. Même si minimes, c'est simple avec ta méthode ou d'autres déjà proposées dans le passé de s'affranchir de ça. Premier point pas génial. Ensuite, pour chaque patient, ils ont donc plusieurs paires de la même acquisition. Ils ne parlent pas du tout de cela. Et ils n'ont que 10 patients pour l'entrainement. MAIS, ils font du 3D et leur réseau ne prend en dimensions d'entrée que des patchs de 32x32x32. Ils extraient 1600 patchs aléatoirement de chaque patient à chaque epoch. Les images faisant 200x200x109, il semble donc y avoir pas mal de redondance dans les patchs. Quoi qu'il en soit, leur méthode semble supposer que N2N n'a qu'à apprendre de manière très local comment le bruit se comporte. La construction de leur base reste largement à discuter. De ton côté, tu utilises toute l'image, certes 2D pour le moment, et chacune de tes paires correspond à une distribution d'activité différente. Difficile de faire mieux car il n'y a aucune rédondance et en incluant bien toute l'image, ton réseau a une chance d'apprendre le fait que la variance ainsi que covariance du bruit est variable spatiallement (ce qui est la réalité). En bref résumé, ne commente pas sur la construction des bases de données, mais mets plutôt l'accent sur le fait qu'ils ont utilisé une simple norme2 dans l'espace image alors que toi tu explores tout un tas de loss en combinant l'espace image et des données. C'est pareil pour l'autre référence où ils proposent d'ajouter une transformation en ondelette, la loss est une unique norm2 entre label et target.}
Both works have limitations since the statistical and physical nature of PET data are under-exploited.

\begin{figure*}[b]
\centering
% trim : left bottom right top 
\begin{minipage}{0.07\linewidth}
    \centering
    \vspace{-10pt}
    \includegraphics[trim={10 0 358 0},clip, height=160pt]{results/cbar.png}

\end{minipage}
\begin{minipage}{0.145\linewidth}
    \begin{subfigure}{\imagewidth}
    \centering
    \vspace{5pt}
    \includegraphics[width=\linewidth,trim={0 20 0 0},clip]{results/gth.png}
    \caption{Ground truth}
    \end{subfigure}
\end{minipage}
%\hfill
\begin{minipage}{0.76\linewidth}
    \centering
\begin{subfigure}{\imagewidth}
    \includegraphics[width=\linewidth,trim={45 45 35 40},clip]{results/prompt_fbp.png}
    \caption{Original FBP}
\end{subfigure}%
% \begin{subfigure}{\imagewidth}  
%     \includegraphics[width=\linewidth,trim={45 45 35 40},clip]{results/nfpt_fbp.png}
%     \caption{Noise-free FBP}
% \end{subfigure}%
\begin{subfigure}{\imagewidth}
    \includegraphics[width=\linewidth,trim={45 45 35 40},clip]{results/image-image_fbp_supervised.png}
    \caption{$\Lc\sym\ImDom$ (supervised)}
\end{subfigure}%
\begin{subfigure}{\imagewidth}
    \includegraphics[width=\linewidth,trim={45 45 35 40},clip]{results/image-image_fbp.png}
    \caption{$\Lc\sym\ImDom$}
\end{subfigure}%
\begin{subfigure}{\imagewidth}
    \includegraphics[width=\linewidth,trim={45 45 35 40},clip]{results/photon-photon_fbp.png}
    \caption{$\Lc\sym\CountDom$}
\end{subfigure}%
\begin{subfigure}{\imagewidth}
    \includegraphics[width=\linewidth,trim={45 45 35 40},clip]{results/photon-image_fbp.png}
    \caption{$\Lc\sym\CI$}
\end{subfigure}%

\begin{subfigure}{\imagewidth}
    \includegraphics[,trim={45 45 35 40},clip,width=\linewidth]{results/prompt_mlem.png}
    \caption{Original MLEM}
\end{subfigure}%
% \begin{subfigure}{\imagewidth}
%     \includegraphics[width=\linewidth,trim={45 45 35 40},clip]{results/nfpt_mlem.png}
%     \caption{Noise-free MLEM}
% \end{subfigure}%
\begin{subfigure}{\imagewidth}
    \includegraphics[width=\linewidth,trim={45 45 35 40},clip]{results/image-image_mlem_supervised.png}
    \caption{$\Lc\sym\ImDom$ (supervised)}
\end{subfigure}%
\begin{subfigure}{\imagewidth}
    \includegraphics[width=\linewidth,trim={45 45 35 40},clip]{results/image-image_mlem.png}
    \caption{$\Lc\sym\ImDom$}
\end{subfigure}%
\begin{subfigure}{\imagewidth}
    \includegraphics[width=\linewidth,trim={45 45 35 40},clip]{results/photon-photon_mlem.png}
    \caption{$\Lc\sym\CountDom$}
\end{subfigure}%
\begin{subfigure}{\imagewidth}
    \includegraphics[width=\linewidth,trim={45 45 35 40},clip]{results/photon-image_mlem.png}
    \caption{$\Lc\sym\CI$}
\end{subfigure}%
\end{minipage}
\caption{Results on the BrainWeb phantom. Images are displayed on the same scale.}
\label{results}
\end{figure*}

%\section{Proposed method}

\begin{comment}
Our first contribution is to make it clear that the N2N principle can be applied without a specific setup involving pairs of sinograms, nor a step consisting in partitioning the recorded events in two sets. Indeed, it is easy to generate two (noisier) sinograms from a single one given that:
\begin{equation*}
%    \begin{cases}
%    Y \sim \text{Poisson}(\lambda) \\
    Y_1 \,|\, Y=y \sim \text{Binomial}(y, 1/2),\qquad
    Y_2 = Y - Y_1,
%    \end{cases}
\end{equation*}
provides two independent variables $Y_1,Y_2 \sim \Pc(\lambda/2)$ if $Y \sim \Pc(\lambda)$.
It is a special case of the Recorrupted2Recorrupted principle~\cite{Monroy24}. 

Secondly, we propose to use this strategy as a common starting point for different denoising frameworks and adapt N2N to PET. We can either apply N2N in the count domain as a pre-processing step before reconstruction, or in the image domain as a post-processing step of an under-regularized reconstruction method. The corresponding loss functions are symmetrized versions $\Lc\sym\CountDom$ and $\Lc\sym\ImDom$ of
%$r$~\eqref{eq:n2n-image}. Let $\hat{x}_i = r(f_\theta(y_i))$ be the reconstructed image.
% introduire les notations dans la phrase, loss function plutot que objective
% \bar{f}(x_1, x_2) pour la symmétrie
%\begin{align}
%    \Lc\CountDom &= \frac{1}{2}\sum_{i=1}^N d\left(f_\theta(y_i^1), y_i^2\right) + d\left(f_\theta(y_i^2), y_i^1\right),
%    \label{eq:n2n-count}
%\\
%\Lc\ImDom &= \frac{1}{2}\sum_{i=1}^N d\left(f_\theta(\hat{x}_i^1), x_i^2\right) + d\left(f_\theta(x_i^2), x_i^1\right).
%    \label{eq:n2n-image}
%\end{align}
\begin{align*}
    \Lc\CountDom(y^1, y^2, \theta) &= \fracp{1}{2}\sum_i%{n=1}^N
\pth{d\left(f_\theta(y_i^1), y_i^2\right) + \mu d\left(f_\theta(y_i^1), y_i^1\right)},
%    \label{eq:n2n-count}
\\
\Lc\ImDom(y^1, y^2, \theta) &= \fracp{1}{2}\sum_i%{i=1}^N
\Bigpth{d\left(f_\theta(\hat{x}_i^1), \hat{x}_i^2\right) + \mu d\Bigpth{\fracp{A}{2}f_\theta(\hat{x}_i^1), y_i^1}}
%    \label{eq:n2n-image}
\end{align*}
where $\hat{x}(y)$ is an under-regularized solution obtained either with FBP or the MLEM algorithm.
Note that in the count domain, $d$ can be chosen as a Kullback-Leibler divergence since in this case, the data is Poisson distributed.

\JI{** ajouter ici une phrase pour parler du 2nd terme des loss.}

When training in the image domain, a physically informed alternative consists in projecting the denoised images back to the count domain and define a symmetrized function $\Lc\sym\CI$ based on
%\begin{equation}
%    \Lc\CI = \fracp{1}{2}\sum_i%{i=1}^N
%    d\left(\fracp{A}{2}f_\theta(x_i^1), y_i^2\right) + d\left(\fracp{A}{2}f_\theta(x_i^2), y_i^1\right).
%    \label{eq:n2n-pi}
%\end{equation}
\begin{equation*}
    \Lc\CI(y^1, y^2, \theta) = \fracp{1}{2}\sum_i%{i=1}^N
    \Bigpth{d\Bigpth{\fracp{A}{2}f_\theta(\hat{x}_i^1), y_i^2} + \mu d\Bigpth{\fracp{A}{2}f_\theta(\hat{x}_i^1), y_i^1}}.
%    \label{eq:n2n-pi}
\end{equation*}
Finally, we add a penalization term to inject some prior knowledge about the raw sinograms to tackle the bias towards smoothness induced by the denoising objectives. The final loss functions are
$\Lc\sym\CountDom + \lambda \Lc\sym\CountDomMC$ and $\Lc\sym\ImDom + \lambda \Lc\sym\ImDomMC$ 
where $\lambda$ is an empirical trade-off parameter, with
\begin{align*}
     \Lc\CountDomMC(y^1, y^2, \theta) &= \sum_i%{i=1}^N
     \Bigpth{d\pth{f_\theta(y_i^1) + f_\theta(y_i^2), y_i^1 + y_i^2}}
%     \label{eq:n2n-i-mc}
     \\
     \Lc\ImDomMC(y^1, y^2, \theta) &= \sum_i%{i=1}^N
     \Bigpth{d\Bigpth{\fracp{A}{4}\left(f_\theta(\hat{x}_i^1) + f_\theta(\hat{x}_i^2)\right), y_i^1 + y_i^2}}
%     \label{eq:n2n-i-mc}
\end{align*}
\end{comment}

\section{Proposed method}

Our first contribution is to make it clear that the N2N principle can be applied without a specific setup involving pairs of sinograms, nor a step consisting in partitioning the recorded events in two sets. Indeed, it is easy to generate two (noisier) sinograms from a single one given that:
\begin{equation*}
%    \begin{cases}
%    Y \sim \text{Poisson}(\lambda) \\
    Y_1 \,|\, Y=y \sim \text{Binomial}(y, 1/2),\qquad
    Y_2 = Y - Y_1,
%    \end{cases}
\end{equation*}
provides two independent variables $Y_1,Y_2 \sim \Pc(\lambda/2)$ if $Y \sim \Pc(\lambda)$.
It is a special case of the Recorrupted2Recorrupted principle~\cite{Monroy24}. 

Secondly, we propose to use this strategy as a common starting point for different denoising frameworks and adapt N2N to PET. We can either apply N2N in the count domain as a pre-processing step before reconstruction, or in the image domain as a post-processing step of an under-regularized reconstruction method.
When training in the image domain, a physically informed alternative consists in projecting the denoised images to compute the error in the count domain.
The corresponding loss functions are symmetrized versions $\Lc\sym\CountDom$, $\Lc\sym\ImDom$ and $\Lc\sym\CI$ of
%$r$~\eqref{eq:n2n-image}. Let $\hat{x}_i = r(f_\theta(y_i))$ be the reconstructed image.
% introduire les notations dans la phrase, loss function plutot que objective
% \bar{f}(x_1, x_2) pour la symmétrie
%\begin{align}
%    \Lc\CountDom &= \frac{1}{2}\sum_{i=1}^N d\left(f_\theta(y_i^1), y_i^2\right) + d\left(f_\theta(y_i^2), y_i^1\right),
%    \label{eq:n2n-count}
%\\
%\Lc\ImDom &= \frac{1}{2}\sum_{i=1}^N d\left(f_\theta(\hat{x}_i^1), x_i^2\right) + d\left(f_\theta(x_i^2), x_i^1\right).
%    \label{eq:n2n-image}
%\end{align}
\begin{align*}
    \Lc\CountDom(y^1, y^2, \theta) &= \sum_i%{i=1}^N
\pth{d_{\text{c}}\left(f_\theta(y_i^1), y_i^2\right) + \mu d_{\text{c}}\left(f_\theta(y_i^1), y_i^1\right)},
%    \label{eq:n2n-count}
\\
\Lc\ImDom(y^1, y^2, \theta) &= \sum_i%{i=1}^N
\Bigpth{d\left(f_\theta(\hat{x}_i^1), \hat{x}_i^2\right) + \mu d_{\text{c}}\Bigpth{\fracp{A}{2}f_\theta(\hat{x}_i^1), y_i^1}},
%    \label{eq:n2n-image}
\\
\Lc\CI(y^1, y^2, \theta) &= \sum_i%{i=1}^N
\Bigpth{d_{\text{c}}\Bigpth{\fracp{A}{2}f_\theta(\hat{x}_i^1), y_i^2} + \mu d_{\text{c}}\Bigpth{\fracp{A}{2}f_\theta(\hat{x}_i^1), y_i^1}},
\end{align*}
where $\hat{x}(y)$ is an under-regularized solution obtained either with FBP or the MLEM algorithm and $d$ and $d_{\text{c}}$  are respectively the L2 loss and a distance measure in the count domain.
Note that $d_{\text{c}}$ can be chosen as a Kullback-Leibler divergence since in this case, the data is Poisson distributed.
The second term is a measure consistency term to ensure that the solution remains true to the data  which avoids over-smoothing. The trade-off must be tuned empirically with the $\mu$ parameter.

\section{Numerical Experiments}

%\subsection{Implementation details}

A numerical protocol has been designed to validate our method prior to an application to real clinical PET data.
% SIMON: Il faut vraiment plus détailler pour donner de la crédibilité à ton travail, du coup je propose ceci pour développer un peu plus. Faudrait vérifier que je ne dis pas de bêtises :-)
We synthetically generated 2,000 2D image phantoms, each of which is an object of uniform activity containing several smaller uniform objects with various hot and cold contrasts.
Sizes, shapes and activity concentrations were all randomly drawn from distributions empirically designed to cover a large spectrum of possible situations.
Images were sampled into $300 \times 300$ voxels of $2 \times 2$ mm$^2$.
2D PET sinograms of 1 million counts were simulated analytically using a parallel-view projector from the \textit{parallelproj} library~\cite{10.3389/fnume.2023.1324562}, with sinogram dimensions of $300 \times 300$ and a radial bin size of 2 mm.
We also included the effect of attenuation, and a resolution model as a 2D image-based and a 1D radial sinogram-based Gaussian filter of 2.8 mm FWHM each.
Random and scatter contributions were also simulated with a fraction of 50\% and 36\%, respectively.
% SIMON: il n'y a plus besoin de dire ça maintenant: Parameters were tuned to match realistic acquisition conditions and scanner geometry.
% SIMON: ancienne version: Here, we synthetically generated image phantoms from which we simulated 1 million count PET sinograms with a parallel-view projector from the \textit{parallelproj} library~\cite{10.3389/fnume.2023.1324562}. We also added the effect of attenuation, and an image- \JI{***or?} a sinogram-based Gaussian filter as a resolution model. Random and scatter events were also simulated with a fraction of 50\% and 36\%, respectively.
%Parameters were tuned to match realistic acquisition conditions and scanner geometry.
This database of 2,000 datasets was used to train three N2N models corresponding to the losses presented in the previous section.
For comparison, we also trained an additional model with loss $\Lc\sym\ImDom$ in a supervised way.
The FBP algorithm was used without regularization and the MLEM algorithm was run for 100 iterations.
%We used a UNet model with 4 million parameters with residual connections within convolutionnal blocks to improve denoising performances.
%FBP and converged MLEM algorithms were used as under-regularized reconstruction methods.
For assessment, we used 2D FDG brain simulated data based on the BrainWeb phantom.
The simulation was done with the same settings as described above.
The performance of the methods were assessed using PSNR and SSIM metrics.

%\subsection{Results}

Figure \ref{results} shows our model predictions on the BrainWeb simulation.
Quantitatively, optimizing either
$\Lc\sym\CountDom$, $\Lc\sym\ImDom$, or $\Lc\sym\CI$
%~\eqref{eq:n2n-count},~\eqref{eq:n2n-image} or~\eqref{eq:n2n-pi}
produced better quality in PSNR and SSIM than raw reconstructions.
As expected, FBP reconstructions (first row) yielded poorer results than MLEM-based ones (second row).
While our three training frameworks give similar results in the MLEM framework, their FBP counterpart gave satisfying results in the case of $\Lc\sym\ImDom$ only.
We also notice that performances are close when comparing self-supervised results and their supervised counterpart for $\Lc\sym\ImDom$, which confirms that the N2N framework is working as intended.
%SIMON: ce n'est plus nécessaire de dire cela: On Figure~\ref{results}, this comparison is only shown when optimizing $\Lc\sym\ImDom$ (column~3 versus column~4).

\section{Conclusion and perspectives}

We demonstrated the feasibility and the effectiveness of a self-supervised deep learning approach to learn to denoise %and/or to reconstruct
from noisy measurements only.%, without ever needing a clean high-count reference data.

% SIMON: j'ajoute une phrase sur le fait que ces résultats sont préliminaires et qu'il faudra mener une étude plus complète avant d'aller plus loin vers le plug et autre
Assessing the proposed self-supervised methods with more simulations and a careful quantitative analysis is required to confirm the preliminary results presented here.
Future work include considering to plug the proposed self-supervised models as a regularization prior into more sophisticated reconstruction algorithms.
%Such a self-supervised model could be plugged as a regularization prior into more sophisticated reconstruction algorithms.
More powerful deep learning approaches for reconstruction, such as unrolling, could also be explored.

%\bibliographystyle{ieeeji}
\bibliographystyle{IEEEtran}
%\bibliography{references}

\begin{thebibliography}{1}

\bibitem{Nuyts02}
{Nuyts {\it et al.}},
%\newblock '',
\newblock {\em {IEEE Trans. Nuclear Sciences}}, vol. 49, no. 1, pp. 56--60,
  2002.

\bibitem{Hashimoto2024}
{Hashimoto {\it et al.}},
%\newblock '',
\newblock {\em {Radiol. Phys. Technol.}}, vol. 17, no. 1, pp. 24–46, 2024.

\bibitem{Gong2019}
{Gong {\it et al.}},
%\newblock '',
\newblock {\em {IEEE Trans. Med. Imaging}}, vol. 38, no. 7, pp. 1655–1665,
  2019.

\bibitem{Lehtinen18}
{Lehtinen {\it et al.}},
%\newblock '',
\newblock in {\em Proc. 35th Int. Conf. Machine Learning}, 2018, vol.~80, pp.
  2965--2974.

\bibitem{Kang2021}
{Kang {\it et al.}},
%\newblock '',
\newblock {\em Electronics}, vol. 10, no. 13, pp. 1529, 2021.

\bibitem{Yie2020}
{Monroy {\it et al.}},
%\newblock '',
\newblock {\em {Nucl. Med. Mol. Imaging}}, vol. 54, no. 6, pp. 299–304, 2020.

\bibitem{Monroy24}
{Monroy {\it et al.}},
%\newblock '',
\newblock Arxiv preprint 2412.04648, 2024.

\bibitem{10.3389/fnume.2023.1324562}
{Schramm {\it et al.}},
%\newblock '',
\newblock {\em Front. Nucl. Med.}, vol. 3, 2024.

\end{thebibliography}

\end{document}
